% SGA 3, Expose III, Loop 1 English reconstruction.
% Scope: 4.5.0 through the 4.7 suppression notice, stopping immediately
% before Remark 4.8.0.
% Authority: Polo--Gille Exp3-14oct24.pdf, component-local pp. 47--51
% (combined-reader pp. 157--161; original side-pages 131--132), with only the
% 4.5.0--4.7 portions of the first and last pages admitted. Force-OCR is a
% locator only. English wording was compared with Jacob C. Reinhold's CC BY
% 4.0 Markdown lineage at commit e7a259f3f8608ad3edf9bf6eead3fd504dd2d23e.
% The Polo--Gille PDF controls prose, mathematics, numbering, notes, and the
% native TeX bicomplex diagram linked below.

\setcounter{footnote}{80}

% Source page: Exp3 local p. 47; combined p. 157; original side-pages 131--132.
\medskip
\noindent\SGAInlineRefTarget{sga3:iii:paragraph:4-5-0}\textbf{4.5.0.}
\footnote{编者注: 新增编号 \SGARefLink{sga3:iii:paragraph:4-5-0}{4.5.0},
以标明此处返回原文.}
沿用 \SGARefLink{sga3:iii:corollary:4-3-1}{4.3.1} 和
\SGARefLink{sga3:iii:corollary:4-4}{4.4} 的记号, 并作如下变换.
拟凝聚 \(\mathscr O_{X_{\mathscr J}}\)-模
\[
 \mathscr I_{Y_{\mathscr J}}/\mathscr I_{Y_{\mathscr J}}^2
\]
被 \(\mathscr I_{Y_{\mathscr J}}\) 零化,
因而是某个拟凝聚 \(\mathscr O_{Y_{\mathscr J}}\)-模的正像.
将此模记作 \(\mathscr N_{Y_{\mathscr J}/X_{\mathscr J}}\),
称为 \(Y_{\mathscr J}\) 在 \(X_{\mathscr J}\) 中的
\emph{余法层 (conormal sheaf)}.\footnote{编者注: 下一句已作修正.}
由于 \(\mathscr J\otimes_{\mathscr O_{S_{\mathscr J}}}
\mathscr O_{Y_{\mathscr J}}\) 被 \(\mathscr I_{Y_{\mathscr J}}\) 零化,
\SGARefLink{sga3:iii:corollary:4-3-1}{4.3.1} 中的交换群层
\(\mathscr A_{\mathscr J}\) 等同于
\[
 \begin{split}
 &\operatorname{Hom}_{\mathscr O_{Y_{\mathscr J}}}
 \left(
  \mathscr I_{Y_{\mathscr J}}/\mathscr I_{Y_{\mathscr J}}^2,
  \mathscr J\otimes_{\mathscr O_{S_{\mathscr J}}}
   \mathscr O_{Y_{\mathscr J}}
 \right) \\
 &\qquad={}
 \operatorname{Hom}_{\mathscr O_{Y_{\mathscr J}}}
 \left(
  \mathscr N_{Y_{\mathscr J}/X_{\mathscr J}},
  \mathscr J\otimes_{\mathscr O_{S_{\mathscr J}}}
   \mathscr O_{Y_{\mathscr J}}
 \right),
 \end{split}
\]
因而对每个 \(i\geq0\),
\[
 H^i(X_{\mathscr J},\mathscr A_{\mathscr J})
 =H^i\!\left(
  Y_{\mathscr J},
  \operatorname{Hom}_{\mathscr O_{Y_{\mathscr J}}}
  \left(
   \mathscr N_{Y_{\mathscr J}/X_{\mathscr J}},
   \mathscr J\otimes_{\mathscr O_{S_{\mathscr J}}}
    \mathscr O_{Y_{\mathscr J}}
  \right)
 \right).
\]

\footnote{编者注: 以下讨论作了扩充.}
现在可按如下方式去掉 \(Y\) 为闭的假设.
首先注意, \(X_{\mathscr J}\) 的每个开子集 \(U_{\mathscr J}\),
都由 \(X\) 的那个与 \(U_{\mathscr J}\) 有相同底拓扑空间的开子概形
\(U\) 经基变换得到.
现在设 \(Y_{\mathscr J}\) 是 \(U_{\mathscr J}\) 的一个闭子概形,
在 \(S_{\mathscr J}\) 上平坦,
以 \(\mathscr I_{Y_{\mathscr J}}\) 表示定义 \(Y_{\mathscr J}\)
的 \(\mathscr O_{U_{\mathscr J}}\) 的拟凝聚理想.
若 \(Y_{\mathscr J}\) 提升为 \(X\) 的一个子概形 \(Y\),
则因 \(Y\) 与 \(Y_{\mathscr J}\) 有相同的底拓扑空间,
它是 \(U\) 的一个闭子概形.
所以, 将 \(Y_{\mathscr J}\) 提升为 \(X\) 或 \(U\)
的一个在 \(S\) 上平坦的子概形, 所遇到的障碍相同; 它位于
\[
 H^1\!\left(
  Y_{\mathscr J},
  \operatorname{Hom}_{\mathscr O_{Y_{\mathscr J}}}
  \left(
   \mathscr N_{Y_{\mathscr J}/X_{\mathscr J}},
   \mathscr J\otimes_{\mathscr O_{S_{\mathscr J}}}
    \mathscr O_{Y_{\mathscr J}}
  \right)
 \right).
\]

最后, 回到 \S0 的记号.
设 \(\mathscr I\) 是 \(\mathscr O_S\) 的一个拟凝聚理想,
满足 \(\mathscr J\subset\mathscr I\) 和 \(\mathscr I\mathscr J=0\),
并设 \(S_0\) 是 \(S_{\mathscr J}\) 的由 \(\mathscr I\) 定义的闭子概形.
对每个 \(S\)-概形 \(Z\), 记
\[
 Z_{\mathscr J}=Z\times_S S_{\mathscr J},
 \qquad Z_0=Z\times_S S_0.
\]
由于 \(\mathscr J\) 被 \(\mathscr I\) 零化,
按 \SGARefLink{sga3:iii:corollary:4-4}{4.4} 的记号有
\[
 \mathscr J\otimes_{\mathscr O_{S_{\mathscr J}}}
  \mathscr O_{Y_{\mathscr J}}
 =\mathscr J\otimes_{\mathscr O_{S_0}}\mathscr O_{Y_0},
\]
以及
\[
 \begin{split}
 &\operatorname{Hom}_{\mathscr O_{Y_{\mathscr J}}}
 \left(
  \mathscr N_{Y_{\mathscr J}/X_{\mathscr J}},
  \mathscr J\otimes_{\mathscr O_{S_{\mathscr J}}}
   \mathscr O_{Y_{\mathscr J}}
 \right) \\
 &\qquad={}
 \operatorname{Hom}_{\mathscr O_{Y_0}}
 \left(
  \mathscr N_{Y_{\mathscr J}/X_{\mathscr J}}
   \otimes_{\mathscr O_{Y_{\mathscr J}}}\mathscr O_{Y_0},
  \mathscr J\otimes_{\mathscr O_{S_0}}\mathscr O_{Y_0}
 \right),
 \end{split}
\]
其余各项也类似. 由此得到:

% Source page: Exp3 local p. 48; combined p. 158; original side-page 132.
\medskip
\noindent\SGAInlineRefTarget{sga3:iii:proposition:4-5}\textbf{命题 4.5.}
\textit{设 \(S\) 是一个概形,
\(S_0\) 与 \(S_{\mathscr J}\) 分别为拟凝聚理想
\(\mathscr I\) 和 \(\mathscr J\) 定义的闭子概形,
其中 \(\mathscr I\supset\mathscr J\) 且 \(\mathscr I\mathscr J=0\).
设 \(X\) 是一个 \(S\)-概形,
\(Y_{\mathscr J}\) 是 \(X_{\mathscr J}\) 的一个子概形,
在 \(S_{\mathscr J}\) 上平坦.
定义 \(\mathscr O_{Y_0}\)-模
\[
 \mathscr A_0=
 \operatorname{Hom}_{\mathscr O_{Y_0}}
 \left(
  \mathscr N_{Y_{\mathscr J}/X_{\mathscr J}}
   \otimes_{\mathscr O_{Y_{\mathscr J}}}\mathscr O_{Y_0},
  \mathscr J\otimes_{\mathscr O_{S_0}}\mathscr O_{Y_0}
 \right).
\]
则:}
\begin{enumerate}
\item[\SGAInlineRefTarget{sga3:iii:proposition:4-5:item:i}(i)]
\textit{存在一个约化为 \(Y_{\mathscr J}\) 且在 \(S\) 上平坦的
\(X\) 的子概形 \(Y\), 当且仅当以下两个条件成立:}
  \begin{enumerate}
  \item[\SGAInlineRefTarget{sga3:iii:proposition:4-5:item:a}(a)]
  \textit{这样的 \(Y\) 在 \(X\) 上局部存在.}
  \item[\SGAInlineRefTarget{sga3:iii:proposition:4-5:item:b}(b)]
  \textit{\(R^1\Gamma(Y_0,\mathscr A_0)\) 中的某个障碍消失.}
  \footnote{编者注: 此处 \(R^1\Gamma(Y_0,\mathscr A_0)\)
  表示 \(\mathscr O_{Y_0}\)-模 \(\mathscr A_0\) 的
  ``凝聚'' 上同调群 \(H^1(Y_0,\mathscr A_0)\),
  用以区别于自 \SGARefLink{sga3:iii:paragraph:4-16}{4.16} 起将考虑的
  ``Hochschild'' 上同调群 \(H^i(Y_0,M_0)\),
  其中 \(Y_0\) 是一个 \(S_0\)-群, \(M_0\) 是一个 \(\mathscr O_{S_0}\)-模.}
  \end{enumerate}
\item[\SGAInlineRefTarget{sga3:iii:proposition:4-5:item:ii}(ii)]
\textit{在这些条件下, 满足上述要求的子概形 \(Y\) 的集合,
是阿贝尔群 \(\Gamma(Y_0,\mathscr A_0)\) 下的主齐性集合.}
\end{enumerate}

\medskip
\noindent\SGAInlineRefTarget{sga3:iii:remark:4-5-1}\textbf{注 4.5.1.}
\footnote{编者注: 此注置于此处, 代替原文的
\SGARefLink{sga3:iii:paragraph:4-7}{注 4.7}.}
由 \SGARefLink{sga3:iii:proposition:4-5:item:ii}{4.5(ii)} 可知,
对 \(X\) 的每对子概形\footnote{编者注: 已将 ``闭子概形'' 改为 ``子概形''.}
\((Y,Y')\), 若它们在 \(S\) 上平坦且约化为 \(Y_{\mathscr J}\),
则存在一个 ``偏差 (deviation)''
\[
 \begin{split}
 d(Y',Y)&\in\Gamma(Y_0,\mathscr A_0)\\
 &=\operatorname{Hom}_{\mathscr O_{Y_0}}
 \left(
  \mathscr N_{Y_{\mathscr J}/X_{\mathscr J}}
   \otimes_{\mathscr O_{Y_{\mathscr J}}}\mathscr O_{Y_0},
  \mathscr J\otimes_{\mathscr O_{S_0}}\mathscr O_{Y_0}
 \right).
 \end{split}
\]
按照 \SGARefLink{sga3:iii:complement:4-4-1}{4.4.1} 采用的符号约定,
\(d(Y',Y)\) 对应于 \(\mathscr O_X\)-模态射
\[
 \mathscr I_{Y'}\lhook\joinrel\longrightarrow
 \mathscr O_X\longrightarrow\mathscr O_Y,
\]
它取值于 \(\mathscr J\mathscr O_Y\simeq
\mathscr J\otimes_{\mathscr O_{S_0}}\mathscr O_{Y_0}\),
且通过 \(\mathscr I_{Y'_{\mathscr J}}=\mathscr I_{Y_{\mathscr J}}\) 分解,
再通过 \(\mathscr N_{Y_{\mathscr J}/X_{\mathscr J}}\) 分解.

\medskip
\noindent\SGAInlineRefTarget{sga3:iii:remark:4-6}\textbf{注 4.6.}
\footnote{编者注: 为便于查阅, 保留了原文的\SGARefLink{sga3:iii:remark:4-6}{注 4.6},
其中没有给出 ``局部完全交'' 的定义.
随后给出取自 SGA 6, VII 1.4 的 ``正确'' 定义
(取代 EGA IV\textsubscript{4}, 16.9.2 的定义),
并给出此注所作的三个断言的证明.}
\textit{若 \(X\) 在 \(S\) 上平坦,
且 \(Y_{\mathscr J}\) 是 \(X_{\mathscr J}\) 中的局部完全交,
则条件 \SGARefLink{sga3:iii:proposition:4-5:item:a}{(a)} 总是满足;
而每个在 \(S\) 上平坦并提升 \(Y_{\mathscr J}\) 的 \(Y\),
都是 \(X\) 中的局部完全交.
此外, 若 \(Y_0\) 仿射, 则条件 \SGARefLink{sga3:iii:proposition:4-5:item:b}{(b)} 也满足.}

\medskip
\noindent\SGAInlineRefTarget{sga3:iii:definition:4-6-1}\textbf{定义 4.6.1.}
\textit{(参见 SGA 6, VII 1.1) 设 \(B\) 是交换环,
\(f:E\to B\) 是一个 \(B\)-线性态射,
其中 \(E\) 是秩为有限数 \(d\) 的自由 \(B\)-模, 并令 \(I=f(E)\).
选定 \(E\) 的一组基后, \(f\) 由 \(B\) 的元素组成的一个
\(d\)-元组 \((f_1,\ldots,f_d)\) 给出,
而 \(I\) 是由这些 \(f_i\) 生成的理想.
\emph{Koszul 复形 (Koszul complex)} \(K_\bullet(f)\)
是分次 \(B\)-模 \(\bigwedge_B^\bullet E\),
配备如下次数为 \(-1\) 的微分:}
\[
 x_1\wedge\cdots\wedge x_i
 \longmapsto
 \sum_{j=1}^i(-1)^{j-1}f(x_j)
 x_1\wedge\cdots\wedge\widehat{x_j}\wedge\cdots\wedge x_i.
\]
于是有一个以 \(B/I\) 为 \(-1\) 次项的增广链复形
\[
 \cdots\longrightarrow\bigwedge^2 E\longrightarrow E
 \xrightarrow{f}B\longrightarrow B/I\longrightarrow0,
\]
由于 \(I=f(E)\), 按定义它在 \(0\) 次正合.
若 \(K_\bullet(f)\) 在 \(>0\) 次无同调,
即上述增广复形是 \(C=B/I\) 的一个解消,
则称态射 \(f\) 是\emph{正则的 (regular)}.

此时, SGA 6, VII 1.2(b) 的证明表明,
对 \(n\in\mathbb N\), \(C\)-模 \(I^n/I^{n+1}\) 都是自由的,
其中 \(I/I^2\) 的秩为 \(d\).

% Source page: Exp3 local p. 49; combined p. 159; original side-page 132.
\medskip
\noindent\SGAInlineRefTarget{sga3:iii:definition:4-6-2}\textbf{定义 4.6.2.}
\textit{(参见 SGA 6, VII 1.4) 设 \(X\) 是概形,
\(Y\) 是子概形, \(U\) 是 \(X\) 的一个开子集,
使得 \(Y\) 是 \(U\) 的由拟凝聚理想 \(\mathscr I_Y\) 定义的闭子概形.}

若 \(Y\hookrightarrow X\) 是 SGA 6, VII 1.4 意义下的正则浸入,
则称子概形 \(Y\) 是 \(X\) 中的\emph{局部完全交 (locally complete intersection)};
即对每个 \(y\in Y\), 存在 \(y\) 在 \(U\) 中的一个仿射开邻域 \(V\),
一个有限自由 \(\mathscr O_V\)-模 \(\mathscr E\),
以及一个像为 \(\mathscr I_Y|_V\) 的正则态射
\(f:\mathscr E\to\mathscr O_V\),
使得 \(K_\bullet(f)\) 是 \(\mathscr O_{Y\cap V}\) 的一个解消.

这蕴含浸入 \(Y\hookrightarrow X\) 局部有限呈示,
且由 \SGARefLink{sga3:iii:definition:4-6-1}{4.6.1}, 余法层
\[
 \mathscr N_{Y/X}=\mathscr I_Y/\mathscr I_Y^2
\]
是一个有限局部自由 \(\mathscr O_Y\)-模.

\medskip
\noindent\SGAInlineRefTarget{sga3:iii:lemma:4-6-3}\textbf{引理 4.6.3.}
\footnote{编者注: 为证明\SGARefLink{sga3:iii:remark:4-6}{注 4.6} 中陈述的结果,
新增了\SGARefLink{sga3:iii:lemma:4-6-3}{引理 4.6.3} 和
\SGARefLink{sga3:iii:lemma:4-6-4}{4.6.4},
\SGARefLink{sga3:iii:proposition:4-6-5}{命题 4.6.5},
以及\SGARefLink{sga3:iii:remark:4-6-6}{注 4.6.6}.}
\textit{设 \(A\) 是一个环, \(J\) 是 \(A\) 的一个平方为零的理想,
置 \(\overline A=A/J\).
设 \(B\) 是一个平坦 \(A\)-代数,
\(E\) 是一个有限秩自由 \(B\)-模,
\(f:E\to B\) 是一个 \(B\)-模态射.
假设由 \(f\) 诱导的态射
\[
 g:\overline E=E\otimes_A\overline A
 \longrightarrow\overline B=B\otimes_A\overline A
\]
正则, 且 \(\overline B/g(\overline E)\) 在 \(\overline A\) 上平坦.
则 \(f\) 正则, 且 \(B/f(E)\) 在 \(A\) 上平坦.}

\noindent\textit{证明.}
置
\[
 C=B/f(E),
 \qquad
 \overline C=C\otimes_A\overline A
 =\overline B/g(\overline E).
\]
首先, 各 \(\bigwedge_B^i E\) 都是自由 \(B\)-模,
因而由于 \(B\) 在 \(A\) 上平坦, 它们是平坦 \(A\)-模.
由于
\[
 \bigwedge_B^\bullet E\otimes_A\overline A
 \simeq\bigwedge_{\overline B}^\bullet\overline E,
\]
有复形正合列
\[
 0\longrightarrow
 J\otimes_A\bigwedge_B^\bullet E
 \longrightarrow\bigwedge_B^\bullet E
 \longrightarrow\bigwedge_{\overline B}^\bullet\overline E
 \longrightarrow0.
\]
此外, 由于 \(J^2=0\), 对每个 \(A\)-模 \(M\) 都有
\[
 J\otimes_A M
 \simeq J\otimes_{\overline A}(\overline A\otimes_A M).
\]
将增广箭头画成虚线, 并将 \(E\) 的秩记为 \(d\),
便得到下列各行正合的双复形.

% Native TeX diagram dependency; source locator follows.
% Authority: Exp3 local p. 50; combined reader p. 160; original side-page 132.
\begin{center}
 \SGAThreeNativeInput[0.78\textwidth]
   {figures/exp3/diagram_038_exp3_p150_koszul_bicomplex.tex}
\end{center}

% Source page: Exp3 local p. 50; combined p. 160; original side-page 132.
左、右两列也正合, 因为 \(K_\bullet(g)\) 是 \(\overline C\) 的一个解消,
且 \(\overline C\) 在 \(\overline A\) 上平坦.
因此, 与不带增广的复形正合列相伴的同调长正合列表明,
\(K_\bullet(f)\) 在 \(>0\) 次无同调, 并且在 \(0\) 次给出正合列
\[
 0\longrightarrow J\otimes_{\overline A}\overline C
 \longrightarrow C\longrightarrow\overline C\longrightarrow0.
\]
由 ``平坦性基本判别法'', \(C\) 在 \(A\) 上平坦
(参见 [BAC], \S III.5, 定理 1).
\hfill\(\square\)

\medskip
\noindent\SGAInlineRefTarget{sga3:iii:lemma:4-6-4}\textbf{引理 4.6.4.}\footnotemark[88]
\textit{设 \(A\) 是交换环, \(J\) 是一个幂零理想,
\(N\subset M\) 是 \(A\)-模, 使得 \(M/N\) 在 \(A\) 上平坦.
若 \(x_1,\ldots,x_n\) 是 \(N\) 中的元素,
其像生成 \(N\) 在 \(M/JM\) 中的像 \(\overline N\),
则它们生成 \(N\).}

事实上, 设 \(N'\) 是由这些 \(x_i\) 生成的 \(N\) 的子模,
并置 \(Q=N/N'\). 态射
\(N'\otimes_A(A/J)\to\overline N\) 是满射.
另一方面, 由于 \(M/N\) 在 \(A\) 上平坦, 态射
\(N\otimes_A(A/J)\to\overline N\) 是双射.
因此 \(Q\otimes_A(A/J)=0\), 从而由 ``幂零 Nakayama 引理'' 得到
\(Q=0\):
\[
 Q=JQ=J^2Q=\cdots=0.
\]

现在可以证明以下命题.

\medskip
\noindent\SGAInlineRefTarget{sga3:iii:proposition:4-6-5}\textbf{命题 4.6.5.}\footnotemark[88]
\textit{设 \(S,\mathscr I,\mathscr J\) 及 \(X,Y_{\mathscr J}\)
如 \SGARefLink{sga3:iii:proposition:4-5}{4.5} 中所述.
再假设 \(X\) 在 \(S\) 上平坦,
且 \(Y_{\mathscr J}\) 是 \(X_{\mathscr J}\) 中的局部完全交.}
\begin{enumerate}
\item[\SGAInlineRefTarget{sga3:iii:proposition:4-6-5:item:a}(a)]
\textit{则 \SGARefLink{sga3:iii:proposition:4-5:item:i}{4.5(i)} 中的条件
\SGARefLink{sga3:iii:proposition:4-6-5:item:a}{(a)} 满足;
此外, 每个在 \(S\) 上平坦并提升 \(Y_{\mathscr J}\) 的 \(Y\),
都是 \(X\) 中的局部完全交.}
\item[\SGAInlineRefTarget{sga3:iii:proposition:4-6-5:item:b}(b)]
\textit{若 \(Y_0\) 还仿射,
则上述结果中的条件 \SGARefLink{sga3:iii:proposition:4-6-5:item:b}{(b)} 也满足.}
\end{enumerate}

% Source page: Exp3 local p. 51; combined p. 161; original side-page 132.
\noindent\textit{证明.}
\SGARefLink{sga3:iii:proposition:4-6-5:item:a}{(a)} 的第一个断言由
\SGARefLink{sga3:iii:lemma:4-6-3}{引理 4.6.3} 得到,
第二个断言随后由\SGARefLink{sga3:iii:lemma:4-6-4}{引理 4.6.4} 得到.
此外, 假设蕴含 (参见 \SGARefLink{sga3:iii:definition:4-6-2}{4.6.2})
\(\mathscr N_{Y_{\mathscr J}/X_{\mathscr J}}\) 是一个有限局部自由
\(\mathscr O_{Y_{\mathscr J}}\)-模.
所以 \(\mathscr O_{Y_0}\)-模
\[
 \mathscr A_0=
 \operatorname{Hom}_{\mathscr O_{Y_0}}
 \left(
  \mathscr N_{Y_{\mathscr J}/X_{\mathscr J}}
   \otimes_{\mathscr O_{Y_{\mathscr J}}}\mathscr O_{Y_0},
  \mathscr J\otimes_{\mathscr O_{S_0}}\mathscr O_{Y_0}
 \right)
\]
拟凝聚 (参见 EGA I, 1.3.12),
因而当 \(Y_0\) 仿射时, 有 \(R^1\Gamma(Y_0,\mathscr A_0)=0\).
\hfill\(\square\)

\medskip
\noindent\SGAInlineRefTarget{sga3:iii:remark:4-6-6}\textbf{注 4.6.6.}\footnotemark[88]
以一个例子结束本段. 此例表明, 在\SGARefLink{sga3:iii:lemma:4-6-3}{引理 4.6.3}
的假设下, 生成理想 \(\overline I=g(\overline E)\) 的正则序列
\((\overline g_1,\overline g_2)\), 未必能够提升为 \(B\) 中的正则序列.

设 \(k\) 是一个域, 置 \(\overline A=k[X,Y]\),
以 \(k\varepsilon\) 表示 \(\overline A\)-模
\(\overline A/(X,Y)\), 从而
\[
 P\varepsilon=P(0,0)\varepsilon
 \qquad(P\in\overline A),
\]
并令 \(A=\overline A\oplus k\varepsilon\),
其中 \(J=k\varepsilon\) 是平方为零的理想.
于是 \(A/J=\overline A\).

代数 \(B=A\otimes_k k[Z,T]\) 在 \(A\) 上自由, 因而平坦,
且 \(\overline B=k[X,Y,Z,T]\). 置
\[
 \overline g_1=XZ-YT,
 \qquad
 \overline g_2=XZ-1.
\]
由于 \(\overline g_1\) 不可约,
\(\overline B/(\overline g_1)\) 是整环;
所以 \((\overline g_1,\overline g_2)\) 是 \(\overline B\) 中的一个正则序列,
它生成理想
\[
 \overline I=(XZ-1,YT-1).
\]
因此
\[
 \overline C=\overline B/\overline I
 =k[X,Y,X^{-1},Y^{-1}]
 =\overline A[X^{-1},Y^{-1}]
\]
是平坦 \(\overline A\)-代数, 也是平坦 \(A\)-代数.
但是, \(\overline g_1\) 到 \(B\) 的每个提升都形如
\[
 XZ-YT+\lambda\varepsilon,
 \qquad \lambda\in k[Z,T],
\]
因而都零化 \(\varepsilon\).

\smallskip
\noindent{\footnotesize\textit{来源说明.}
Polo--Gille PDF 将此提升印为 \(XY-ZT+\lambda\varepsilon\).
由于所提升的元素是 \(\overline g_1=XZ-YT\),
上述提升公式在加上 \(\lambda\varepsilon\) 之前恢复了同一个多项式.}

\medskip
\noindent\SGAInlineRefTarget{sga3:iii:paragraph:4-7}\textbf{4.7.}
\SGARefLink{sga3:iii:paragraph:4-7}{注 4.7} 已从此处移去,
置于 \SGARefLink{sga3:iii:remark:4-5-1}{4.5.1}.

% Continue with Remark 4.8.0 on Exp3 local p. 51 / combined p. 161.

% Translation coverage: all of 4.5.0 through the original 4.7 relocation notice.
% Statistics: source lines=447; target lines=456; inline math=246; display math=26; equation/align/multline blocks=0; unresolved translation TODOs=0.
